\section{Augmentation, normal products, and the subregular Poisson quotient}
\label{sec:cbc28-subregular-bar}

The strict reduction map has the affine--ghost complex as its target.
The vacuum product $J_{(1)}J$ prevents replacing that target by the bare affine algebra at level zero.
It also restricts augmented bar constructions before any question of completion arises.

Use the universal polynomial family $W_R\subset C_R$ of
Theorem~\ref{thm:cbc27-subregular-map}, with $R=\mathbb C[k]$.
The fields $J,U,P,V$ retain their definitions and cohomological degree zero.
At a complex level $k$, write $W_k,C_k$ for their fibers.
A vertex augmentation is a unital vertex homomorphism to $\mathbb C$ with zero translation and only its ordinary multiplication.

\begin{proposition}[The affine target and the vacuum term]\label{prop:cbc28-nonembedding}
There is no unital vertex homomorphism
\[
 W_0\longrightarrow V^0(\mathfrak{sl}_3).
\]
More generally, $W_k$ has a vertex augmentation if and only if $k=-3/2$.
No pair $W_k,W_{-k-6}$ admits vertex augmentations on both factors.
\end{proposition}
\begin{proof}
The universal affine algebra $V^0(\mathfrak{sl}_3)$ has an augmentation that kills every current.
Its defining brackets permit this map: the affine central term is zero at level zero, and every Lie bracket maps to zero.
Composing a proposed map from $W_0$ with this augmentation gives a vertex augmentation of $W_0$.
However, $J_{(1)}J=\mathbf1$ in $W_0$, while every positive product in $\mathbb C$ vanishes.
This would give $0=1$.

At arbitrary $k$, the same argument applied to
$J_{(1)}J=(2k+3)\mathbf1/3$ requires $k=-3/2$.
At that level, all scalar terms in the universal presentation
\eqref{eq:cbc27-j-products}--\eqref{eq:cbc27-g-products} vanish.
Sending $J,G^+,G^-,T$ to zero therefore defines a vertex augmentation by the presentation.
Here $k+3=3/2$, so this presentation is defined.
Finally,
\[
 \frac{2k+3}{3}+\frac{2(-k-6)+3}{3}=-2.
\]
Both coefficients cannot vanish over $\mathbb C$.
The same contradiction applies over every nonzero $\mathbb C$-algebra.
\end{proof}

The ghost target has no such augmentation because $b_i{}_{(0)}c_i=\mathbf1$.
Thus the explicit map $W_0\hookrightarrow C_0$ is consistent with this obstruction.
The proposition excludes a bare-affine inclusion, not the reduction map.

\begin{proposition}[The first normal-product associator]\label{prop:cbc28-bar-obstruction}
At every complex level, including $k=-3$, the polynomial fields satisfy
\begin{equation}\label{eq:cbc28-associator}
 :(:JJ:)U:-:J(:JU:):=2:(\partial J)U:\ne0.
\end{equation}
Consequently the ordinary tensor-bar coderivation defined only from the normal product does not square to zero.
This failure occurs at tensor length three.
There is no uncurved $A_\infty$ structure with $m_1=0$ and $m_2(a,b)=:ab:$ on this state space.
\end{proposition}
\begin{proof}
For even states $a,b$, the vertex iterate identity gives
\[
 (:ab:)_{(-1)}c-a_{(-1)}b_{(-1)}c
 =\sum_{r\geq0}\left(a_{(-r-2)}b_{(r)}c+b_{(-r-2)}a_{(r)}c\right).
\]
Apply this with $a=b=J$ and $c=U$.
The relations $J_{(0)}U=U$ and $J_{(r)}U=0$ for $r\geq1$ leave exactly $2J_{(-2)}U$.
Translation identifies this state with $2:(\partial J)U:$.
Its leading PBW symbol is the nonzero monomial $2(\partial J)U$ in the polynomial jet algebra.
Theorem~\ref{thm:cbc27-subregular-map} proves that it remains nonzero in every fiber.

Use the cohomological desuspension $sF=F[1]$, with $|sa|=|a|-1$.
On $T^c(sF)=\bigoplus_{n\geq0}(sF)^{\otimes n}$, let $b_2$ be the coderivation with
$b_2(sa\otimes sb)=(-1)^{|a|}s(:ab:)$.
For these degree-zero fields, its square gives
\[
 b_2^2(sJ\otimes sJ\otimes sU)=2s\bigl(:(\partial J)U:\bigr)\ne0.
\]
At $k=-3/2$, the same word lies in the augmentation ideal.
Thus passing to that ideal does not remove this failure.
In an uncurved $A_\infty$ algebra with $m_1=0$, the arity-three identity is associativity of $m_2$.
Higher $m_n$ cannot change that identity.
\end{proof}

The proposition concerns this specified normal-product bar.
A curved chiral bar has additional operations and configuration-space coefficients.
Its definition must supply the terms that replace the failed identity~\eqref{eq:cbc28-associator}.
The obstruction does not exclude a separately constructed curved comparison at reflected levels.

\subsection{The finite Poisson reduction}

For a differential vertex superalgebra $F$, define
\[
 C_2(F)=\operatorname{span}\{a_{(-2)}b:a,b\in F\},\qquad
 R_F=F/C_2(F).
\]
The differential preserves $C_2(F)$.
The quotient has an associative supercommutative product and a degree-zero super-Poisson bracket
\[
 \bar a\,\bar b=\overline{a_{(-1)}b},\qquad
 \{\bar a,\bar b\}=\overline{a_{(0)}b}.
\]
Translation vanishes in this quotient.
Quasi-commutativity, the iterate identity, and the vertex Jacobi identity prove these assertions.
They are the $C_2$ construction in \cite{AffineDSArakawa026}*{Proposition~3.15, p.~21}.
The quotient map retains these two induced operations.
It is not a vertex homomorphism to a constant-field vertex algebra.

The affine and charged-fermion PBW bases identify
\begin{equation}\label{eq:cbc28-classical-carrier}
 R_{C_R}=K_R
 =R[a,e,h,j,u,v,A,B]\otimes\Lambda(b_1,b_2,c_1,c_2).
\end{equation}
Here the same symbols denote their quotient classes.
The bracket on the affine variables is matrix commutation, while
$\{b_i,c_j\}=\{c_j,b_i\}=\delta_{ij}$.
The affine central coefficient and every derivative have disappeared.
For the affine quotient this is \cite{AffineDSArakawa026}*{Eq.~(3.19), p.~24}.
For the charged ghosts it is the calculation in \cite{AffineDSArakawa026}*{Example~4.10, pp.~27--28}.
Their tensor product gives~\eqref{eq:cbc28-classical-carrier}.
The same ordered-mode argument applies over $R$ because all defining coefficients are polynomial in $k$.

Use ordinary supercommutative multiplication in $K_R$ and set
\[
 H=h+2b_1c_1+b_2c_2,\quad J=j+b_2c_2,\quad
 U=u-b_2c_1,\quad X=v-b_1c_2,
\]
\begin{equation}\label{eq:cbc28-finite-invariants}
 p=A+H^2/4+UX,\qquad v_0=B-HX/2-3JX/2.
\end{equation}
Let $A_R=R[J,U,p,v_0]$, with zero differential.
The inclusion $i:A_R\to K_R$ is given by these explicit polynomials.

\begin{theorem}[Reduction after the $C_2$ quotient]\label{thm:cbc28-c2-reduction}
The natural map induced by $W_R\hookrightarrow C_R$ identifies $R_{W_R}$ with $A_R$ and gives a quasi-isomorphism
\[
 (R_{W_R},0)=(A_R,0)\xrightarrow{\ i\ }(K_R,D)
\]
of differential Poisson superalgebras over $R$.
Its cohomology is concentrated in degree zero, and these assertions commute with every specialization of $k$.
The Poisson bracket on $A_R$ is
\begin{align*}
 \{J,U\}&=U,&\{J,v_0\}&=-v_0,&\{J,p\}&=0,\\
 \{U,v_0\}&=p-\tfrac94J^2,&
 \{p,U\}&=-\tfrac32JU,&\{p,v_0\}&=\tfrac32Jv_0.
\end{align*}
\end{theorem}
\begin{proof}
The differential on $K_R$ is the differential of $C_R$ with every derivative set to zero.
Write $a'=a-1$.
The triangular polynomial change to coordinates
\[
 (a',e,b_1,b_2,H,X,c_1,c_2,J,U,p,v_0)
\]
is invertible: first recover $h,j,u,v$ from $H,J,U,X$ and the ghosts, then recover $A,B$ from~\eqref{eq:cbc28-finite-invariants}.
In these coordinates,
\[
 D=a'\partial_{b_1}+e\partial_{b_2}-2c_1\partial_H+c_2\partial_X,
 \qquad DJ=DU=Dp=Dv_0=0.
\]
Odd derivatives act from the left.
The odd operator
\[
 s=b_1\partial_{a'}+b_2\partial_e-\tfrac12H\partial_{c_1}+X\partial_{c_2}
\]
satisfies $Ds+sD=N$, where $N$ counts the first eight coordinates.
Division by the positive integer eigenvalues of $N$ contracts their positive-degree part.
The projection $\pi:K_R\to A_R$ sets those eight coordinates to zero.
It is an algebra chain map, $\pi i=1$, and $Ds/N+sD/N=1-i\pi$ on the positive-$N$ part.
Thus $i$ is a quasi-isomorphism of differential associative algebras.

Every class in $R_{W_R}$ is a polynomial in the images of the four strong generators.
Indeed, a derivative belongs to $C_2(W_R)$, and the quotient product is associative.
The resulting surjection $A_R\to R_{W_R}$ composed with
$R_{W_R}\to H^0(K_R)$ is the isomorphism just proved.
It follows that the surjection is injective and that $R_{W_R}=A_R$.
Functoriality of $C_2$ makes $i$ a Poisson map.
Taking the zero products of the polynomial fields, or substituting their displayed polynomials in the matrix and ghost brackets, gives the six formulas.
No inverse of a polynomial in $k$ occurs.
The coordinate change and contraction therefore prove the same assertions in every fiber.
\end{proof}

The projection $\pi$ need not preserve the Poisson bracket.
For example, $\pi h=0$, $\pi a=1$, and
$\pi\{h,a\}=2\ne\{\pi h,\pi a\}$.
This distinction matters for the bar construction below, which uses associative multiplication.

\subsection{An actual associative bar comparison}

Give $A_R$ the augmentation $\epsilon_A(J)=\epsilon_A(U)=\epsilon_A(p)=\epsilon_A(v_0)=0$.
Give $K_R$ the chain augmentation $\epsilon_K=\epsilon_A\pi$.
In original coordinates it sends $a$ to one and every other affine or ghost variable to zero.
These are augmentations of differential associative algebras.
They are not asserted to preserve the Poisson bracket or lift to vertex augmentations.
Both $i$ and $\pi$ preserve them.

For an augmented differential associative algebra $E$, put $\bar E=\ker\epsilon_E$ and
\[
 B(E)=\bigoplus_{n\geq0}(s\bar E)^{\otimes_R n},\qquad |sa|=|a|-1.
\]
Use deconcatenation and the tensor Koszul signs.
Its bar differential is the coderivation determined by
\[
 b_1(sa)=-s(Da),\qquad
 b_2(sa\otimes sb)=(-1)^{|a|}s(ab).
\]

\begin{theorem}[Direct-sum bar comparison]\label{thm:cbc28-bar-comparison}
The maps $B(i)$ and $B(\pi)$ are explicit coalgebra chain maps and quasi-isomorphisms
\[
 B(A_R)\underset{B(\pi)}{\overset{B(i)}{\rightleftarrows}}B(K_R),
 \qquad B(\pi)B(i)=1.
\]
On a word, they apply $i$ or $\pi$ separately to every letter.
The construction commutes with every specialization of $k$.
\end{theorem}
\begin{proof}
The differential squares to zero by $D^2=0$, the Leibniz rule, and associativity.
The signs displayed above give these three identities in arities one, two, and three, respectively.
The maps preserve the unit, augmentation, multiplication, and differential, so their tensorwise extensions commute with $b_1,b_2$ and deconcatenation.

The contraction in Theorem~\ref{thm:cbc28-c2-reduction} preserves the augmentation kernels.
All these complexes and maps are obtained from complexes of $\mathbb C$-vector spaces by tensoring with $R$.
It follows that every finite tensor power of $s\bar i$ is a quasi-isomorphism.
Filter $B(i)$ by tensor length.
The quotient at length $n$ is that tensor-power map with its internal differential.
Its cone is acyclic.
For any cycle in the total cone, solve the boundary equation in its highest length quotient and subtract that boundary.
The largest nonzero length decreases.
Since every word sum is finite, this procedure terminates and proves that $B(i)$ is a quasi-isomorphism.
The equation $B(\pi)B(i)=1$ and the same argument prove the assertion for $B(\pi)$.
The contraction works over every fiber, so specialization gives the same maps and proof.
\end{proof}

This bar comparison applies to the full four-generator $C_2$ quotient of the subregular algebra.
The quotient construction and Theorem~\ref{thm:cbc28-c2-reduction} provide its comparison with the affine--ghost quotient.
It retains no first-product operation and kills every derivative and the associator~\eqref{eq:cbc28-associator}.
Indeed, $\overline{\partial J}=0$ but
$\overline{(\partial J)_{(1)}U}=-U\ne0$.
Thus the first product cannot descend to this quotient at any level.
This prevents using the theorem as a proof about the unreduced vertex carrier or a completed chiral bar.

\begin{corollary}\label{cor:cbc28-bar-cohomology}
Let $z_1=J,z_2=U,z_3=p,z_4=v_0$ in $A_R$.
The cohomology of either bar in Theorem~\ref{thm:cbc28-bar-comparison} is
\[
 H^{-n}\simeq\Lambda_R^n(R^4)\quad(0\leq n\leq4),
 \qquad H^r=0\quad\text{otherwise}.
\]
The class indexed by $i_1<\cdots<i_n$ is represented in $B(A_R)$ by
\[
 \sum_{\sigma\in S_n}\operatorname{sgn}(\sigma)
 [z_{i_{\sigma(1)}}|\cdots|z_{i_{\sigma(n)}}].
\]
Applying $B(i)$ gives its explicit representative in $B(K_R)$.
\end{corollary}
\begin{proof}
The normalized free bar resolution has terms
$A_R\otimes_R\bar A_R^{\otimes n}$ in degree $-n$ and differential
\[
 \begin{split}
 d(a_0|a_1|\cdots|a_n)
 ={}&-a_0a_1|a_2|\cdots|a_n\\
 &+\sum_{i=1}^{n-1}(-1)^{i-1}
 a_0|\cdots|a_ia_{i+1}|\cdots|a_n.
 \end{split}
\]
Its augmentation is $\epsilon_A$.
The operator $h(a_0|a_1|\cdots|a_n)=-1|\bar a_0|a_1|\cdots|a_n$, with
$\bar a_0=a_0-\epsilon_A(a_0)$, satisfies $dh+hd=1-\eta\epsilon_A$.
Here $\eta:R\to A_R$ is the unit in degree zero, and the last term is zero in negative degrees.
Thus this is a free resolution of $R$ over $A_R$.

The polynomial Koszul complex $A_R\otimes\Lambda(\eta_1,\ldots,\eta_4)$, with
$|\eta_i|=-1$ and $d\eta_i=-z_i$, is another free resolution.
The sequence $z_1,\ldots,z_4$ is regular: each successive quotient is a polynomial ring in the remaining variables over $R$.
Tensoring its four two-term resolutions proves exactness in negative degrees.
The antisymmetrization in the statement, with initial coefficient $1$, defines a chain map from this Koszul resolution to the bar resolution.
Its internal product terms cancel in pairs by commutativity, while its first terms reproduce $d\eta_i=-z_i$.
It lifts the identity of $R$.
The usual inductive lifting of maps between projective resolutions gives a homotopy inverse.
Tensoring with $R$ over $A_R$ therefore computes the bar cohomology by the zero-differential exterior algebra on the four $\eta_i$.
Theorem~\ref{thm:cbc28-bar-comparison} supplies the same result for $B(K_R)$.
\end{proof}

\begin{theorem}[Induced-module bar comparison]\label{thm:cbc28-module-bar}
Let $M$ be a differential left $A_R$-module and put
$M_K=K_R\otimes_{A_R}M$ using $i$.
For $E=A_R,K_R$, form the direct-sum module bar
\[
 B(E;M_E)=\bigoplus_{n\geq0}(s\bar E)^{\otimes_R n}\otimes_R M_E,
 \qquad M_{A_R}=M.
\]
Its differential consists of the internal differentials, the adjacent $b_2$ above, and the last-letter action
$sa\otimes m\mapsto am$, with the tensor Koszul signs.
The maps
\[
 [a_1|\cdots|a_n]\otimes m
 \longmapsto [i(a_1)|\cdots|i(a_n)]\otimes(1\otimes m)
\]
and the reverse map applying $\pi$ to the letters and $x\otimes m\mapsto\pi(x)m$ to the module are chain quasi-isomorphisms.
Their reverse composite on $B(A_R;M)$ is the identity.
They commute with the natural bar comodule structures relative to $B(i)$ and $B(\pi)$.
\end{theorem}
\begin{proof}
The module differential squares to zero by the differential-module Leibniz identity and the associative action identity.
For example, on two degree-zero letters its action terms are
$(ab)m-a(bm)=0$; on a homogeneous last letter the sign-free action anticommutes with the internal differential because $b_1(sa)=-s(Da)$.
The displayed maps preserve all these operations.

In the coordinates of Theorem~\ref{thm:cbc28-c2-reduction}, $K_R=A_R\otimes_R Q_R$, where $Q_R$ is the algebra on the eight contracted coordinates.
The contraction $s/N$ is $A_R$-linear.
It therefore gives a contraction of $M_K\rightleftarrows M$ for every differential $M$, with the tensor signs.
The contraction of $s\bar K_R\rightleftarrows s\bar A_R$ likewise gives homotopy equivalences on all finite tensor powers, including after tensoring with $M$.
The associated map at every tensor length is consequently a quasi-isomorphism.
The finite highest-length argument in Theorem~\ref{thm:cbc28-bar-comparison} proves the assertion on the direct sums.
This proof does not require boundedness of $M$ or completion.
\end{proof}

For $M=R$ with its augmentation action, $B(A_R;R)=B(A_R)$, so the theorem computes the cohomology of $B(K_R;K_R\otimes_{A_R}R)$ by the same exterior algebra.
For a general affine vertex module there is no assertion here that its $C_2$ quotient is an induced module $M_K$.
Constructing that comparison, and restoring the lost first products, are separate requirements.

\subsection{Changing the central-charge parameter}

Over $\mathbb C[k,(k+3)^{-1}]$, the central charge is $c=25-6t-24/t$, where $t=k+3$.
Changing the base from the level to $c$ is the finite algebra
\begin{equation}\label{eq:cbc28-central-base}
 S=\mathbb C[c,t]/\bigl(6t^2+(c-25)t+24\bigr).
\end{equation}
Its element $t$ is invertible, with $t^{-1}=-(6t+c-25)/24$.
Eliminating $c$ gives $S\simeq\mathbb C[t,t^{-1}]$.
As a $\mathbb C[c]$-module, $S$ is free on $1,t$ because the relation is monic after division by $6$.
Its discriminant is $(c-1)(c-49)$.
Thus the two branches are exchanged by $t\mapsto4/t$ and meet at $c=1,49$.
These equations define the coefficient change without a square root of $c$ or an adic completion.
Choosing a formal branch at a specified point is an additional base change.

\subsection{Orbit dimensions and generator weights}

\begin{proposition}\label{prop:cbc28-sl4-centralizers}
For the Dynkin good grading of each nilpotent orbit in $\mathfrak{sl}_4$, the universal reduction at $k\ne-4$ has the following even strong generators.
\[
 \begin{array}{c|c|c|c}
 \text{partition}&\dim\mathcal O&\dim\mathfrak{sl}_4^f&\text{generator weights}\\\hline
 (1,1,1,1)&0&15&1^{15}\\
 (2,1,1)&6&9&1^4,(3/2)^4,2\\
 (2,2)&8&7&1^3,2^4\\
 (3,1)&10&5&1,2^3,3\\
 (4)&12&3&2,3,4
 \end{array}
\]
An exponent records multiplicity.
In particular, the subregular $\mathfrak{sl}_4$ algebra has five generators, with integral weights $1,2,2,2,3$.
\end{proposition}
\begin{proof}
Let $L_m$ be the irreducible $\mathfrak{sl}_2$-module of highest weight $m$ and dimension $m+1$.
A Jordan block of size $d$ gives the summand $L_{d-1}$ in the defining representation.
The decompositions of its traceless endomorphisms, in the order of the table, are
\[
 15L_0,\qquad 4L_0\oplus4L_1\oplus L_2,\qquad
 3L_0\oplus4L_2,\qquad L_0\oplus3L_2\oplus L_4,
 \qquad L_2\oplus L_4\oplus L_6.
\]
They follow from
$L_a\otimes L_b=\bigoplus_{r=0}^{\min(a,b)}L_{a+b-2r}$ and removal of the scalar endomorphism.
In each $L_m$, the kernel of the lowering operator is one-dimensional and has grading degree $-m/2$.
Kac--Wakimoto's structure theorem therefore gives one even strong generator of weight $1+m/2$ for each summand
\cite{SubregularKW027}*{Theorem~4.1, pp.~10--12}.
Counting these summands gives the centralizer dimensions.
Subtracting each from $\dim\mathfrak{sl}_4=15$ gives the orbit dimensions.

For the subregular row, this can also be checked directly.
Take $f=E_{21}+E_{32}$ and $x=\operatorname{diag}(1,0,-1,0)$.
Its centralizer has basis
\[
 \operatorname{diag}(1,1,1,-3),\quad f,\quad E_{34},\quad E_{41},\quad f^2,
\]
with grading degrees $0,-1,-1,-1,-2$.
\end{proof}

\begin{corollary}\label{cor:cbc28-principal-affine}
For $n\geq2$ and noncritical levels $k,\ell\ne-n$, there is no conformal vertex isomorphism
\[
 \mathcal W^k(\mathfrak{sl}_n,f_{\mathrm{prin}})
 \simeq V^\ell(\mathfrak{sl}_n)
\]
between the universal algebras with their standard conformal vectors.
\end{corollary}
\begin{proof}
The principal centralizer has basis $f,f^2,\ldots,f^{n-1}$ with grading degrees $-1,-2,\ldots,-(n-1)$.
The same structure theorem gives strong generator weights $2,3,\ldots,n$.
Its weight-one space is therefore zero.
The universal affine algebra has weight-one space $\mathfrak{sl}_n$, of dimension $n^2-1$.
A conformal isomorphism preserves these weight spaces.
\end{proof}

Partition transpose describes nilpotent orbit duality in type $A$.
The weight-space obstruction shows that this finite combinatorial operation does not, by itself, identify the vertex algebras attached to the two orbits.

\begin{proposition}[A vacuum obstruction to generator truncation]\label{prop:cbc28-truncation}
Let $\mathcal W_3$ have the normalization $W_{(5)}W=(c/3)\mathbf1$.
If $c\ne0$, no unital vertex homomorphism out of this algebra can send $W$ to zero.
In particular, at $c=-30$, deleting the weight-three generator cannot define a map to a Virasoro algebra.
\end{proposition}
\begin{proof}
Applying such a homomorphism to $W_{(5)}W=(c/3)\mathbf1$ would give $0=(c/3)\mathbf1$ in its nonzero target.
At $c=-30$, the right side is $-10\mathbf1$.
\end{proof}
